\section{Limitations in detail}\label{app:limits}
\begin{itemize}\setlength\itemsep{1pt}
\item \textbf{Hardware.} Power limits differed between machines. Rented machines are often shared or unsteady: of
\jmServAll{} launches on server processors only \jmServWord{} passed every check. In the two tests of the time relation on new
machines, \jlGated{} and \jmGated{} machines stopped at the gate and \jlInvalid{} and \jmInvalid{} ran unsteadily.
The first RTX 4090 test deviated from its registration in three ways: its output check used the RTX 5090's reference
loss, so its machines were relaunched with the RTX 4090's; two relaunches that failed to download the model were
repeated once; and one of the four machines eligible for relaunch was not relaunched. In job 112 one
machine never finished downloading the model, and its replacement started before it was given up. Job 113 ran on two
machines and stopped there, as registered. Job 114 started seven machines, its registered limit: three had GPUs rented
before and one ran unsteadily, so three were valid, not the five it aimed for; its replacements skipped a cheaper listed
offer each time, a deviation from its registered order (\cref{app:prereg}). Job 115 started \dqStartedWord, its limit: \dqGated{}
stopped at a gate, and the two valid machines both have a link-to-CPU ratio of \dqRatioRng.
\item \textbf{The bounds.} MIN's count lets an expert served in a step be evicted later in that step; keeping it for
the step raises the bound slightly (\cref{tab:limit}). $R^\star$ is counted on the 30-problem trace; on the first 20
problems, which the later experiments time, it changes by \rsDiffLow\%. The rate check of \cref{tab:robust} uses the
closed-form account of \cref{app:relation}. $T_{\text{GPU}}$ counts every kernel of an Nsight profile but the
expert GEMVs, the cache's control kernels and copies, and the wait for the CPU helper (activation quantization is in
it); the \slProfN{} profiles are those of jobs 069c and 105 at both gpt-oss budgets, on \slProfHosts{} machines
(\texttt{scripts/sumlaw\_paper.py}). It is not a property of the machine alone; with the RTX
4090's own profile (\cxTgpuLow\,ms) the residual on the slow-link RTX 4090s is \cxResidOwnLowMin--\cxResidOwnLowMax\%
of the gap at 11\%. Measured against \cref{eq:dep}, the tighter bound on slow links, the read-ahead oracle still leaves
\dmEqThreeLeftSlowLowMin--\dmEqThreeLeftSlowLowMax\% of the gap to it on the two slow-link panel machines and
\dmEqThreeLeftCardLowMin--\dmEqThreeLeftCardLowMax\% on the three slow-link RTX 4090s, against a median of
\dmEqThreeLeftFastLowMed\% on fast links: the oracles fetch their misses over the link, the slow path there.
\item \textbf{The decomposition.} The GPU's non-expert time is profiled, not removed. The machines pool runs of 20 and
30 problems and three versions of the engine. \Cref{eq:dep} assumes that the CPU reads an expert's weights when it
runs the expert. The shares depend on the probe through \cref{eq:limit}: with every machine's $B_{\mathrm{host}}$ 10\%
(22\%) higher, the fast-link shares of \cref{tab:gap} at 11\% move by at most \dmProbeTenTogetherLowMax{}
(\dmProbeTwentyTwoTogetherLowMax) points for \emph{Both} and \dmProbeTenLeftLowMax{} (\dmProbeTwentyTwoLeftLowMax) for
\emph{left}. The fetches-off column of \cref{tab:gap} rests on \dpNWord{} machines of two kinds: \dpHiRatioN{} Ryzen 9 5950Xs of one
provider (ratio \dpRatioHiRatioRng) and \dpLoRatioN{} Intel desktop processors (ratio \dpRatioLoRatioRng).
\item \textbf{The oracles.} In the engine the fewest-admission schedule misses \jfMissDevMin--\jfMissDevMax\% more than
the greedy one, more than its replay predicts; its measured gains include that cost. Degraded windows draw errors
independently, while a real predictor's errors correlate with what the cache holds.
\end{itemize}
